\section{Results}

\subsection{Benchmark Composition}

We begin by summarizing the composition of Dataset~A as an empirical artifact rather than only as input to the baseline comparison. Let $N_{\mathit{seed}}$ denote the number of verified seed programs, $N_{\mathit{items}}$ the number of transformed nearby-version instances retained after the generation pipeline, and $N_{\mathit{eval}}$ the number of items that survive to baseline evaluation. The first summary table should report these quantities by seed-program stratum and by refactoring class, together with the ratio $N_{\mathit{eval}} / N_{\mathit{items}}$. This view establishes whether the benchmark has reasonable coverage across both source regimes and transformation types, and it also reveals whether some refactoring families are too sparse for stable table-level reporting in an early experiment tier. The broader pilot in Section~4 suggests that the composition table is informative once multiple strata are present, even before all seven classes are fully populated.

\begin{table}[t]
\centering
\caption{Planned benchmark-composition summary. The table is intended to show that Dataset~A covers multiple seed strata and refactoring classes while retaining enough valid nearby-version pairs for baseline evaluation.}
\label{tab:benchmark-composition}
\begin{tabular}{lrrr}
\toprule
Stratum / class & Seeds & Items & Eval-ready \\
\midrule
Tutorial / educational & \tbdcell{$N^{\mathrm{tut}}_{\mathit{seed}}$} & \tbdcell{$N^{\mathrm{tut}}_{\mathit{items}}$} & \tbdcell{$N^{\mathrm{tut}}_{\mathit{eval}}$} \\
Benchmark / suite & \tbdcell{$N^{\mathrm{bench}}_{\mathit{seed}}$} & \tbdcell{$N^{\mathrm{bench}}_{\mathit{items}}$} & \tbdcell{$N^{\mathrm{bench}}_{\mathit{eval}}$} \\
Proof-engineering-heavy & \tbdcell{$N^{\mathrm{heavy}}_{\mathit{seed}}$} & \tbdcell{$N^{\mathrm{heavy}}_{\mathit{items}}$} & \tbdcell{$N^{\mathrm{heavy}}_{\mathit{eval}}$} \\
\midrule
Rename & \tbdcell{$-$} & \tbdcell{$N^{\mathrm{rename}}_{\mathit{items}}$} & \tbdcell{$N^{\mathrm{rename}}_{\mathit{eval}}$} \\
Helper extraction & \tbdcell{$-$} & \tbdcell{$N^{\mathrm{extract}}_{\mathit{items}}$} & \tbdcell{$N^{\mathrm{extract}}_{\mathit{eval}}$} \\
Inlining & \tbdcell{$-$} & \tbdcell{$N^{\mathrm{inline}}_{\mathit{items}}$} & \tbdcell{$N^{\mathrm{inline}}_{\mathit{eval}}$} \\
Statement reordering & \tbdcell{$-$} & \tbdcell{$N^{\mathrm{reorder}}_{\mathit{items}}$} & \tbdcell{$N^{\mathrm{reorder}}_{\mathit{eval}}$} \\
Loop reshaping & \tbdcell{$-$} & \tbdcell{$N^{\mathrm{loop}}_{\mathit{items}}$} & \tbdcell{$N^{\mathrm{loop}}_{\mathit{eval}}$} \\
Contract factoring & \tbdcell{$-$} & \tbdcell{$N^{\mathrm{contract}}_{\mathit{items}}$} & \tbdcell{$N^{\mathrm{contract}}_{\mathit{eval}}$} \\
Ghost-code movement & \tbdcell{$-$} & \tbdcell{$N^{\mathrm{ghost}}_{\mathit{items}}$} & \tbdcell{$N^{\mathrm{ghost}}_{\mathit{eval}}$} \\
\midrule
Total & \tbdcell{$N_{\mathit{seed}}$} & \tbdcell{$N_{\mathit{items}}$} & \tbdcell{$N_{\mathit{eval}}$} \\
\bottomrule
\end{tabular}
\end{table}


\subsection{Breakage Rate by Refactoring Class}

The core descriptive result is how often verification breaks after semantics-preserving refactoring. For each refactoring class $class_i$, we report the number of attempted transformations, the number of broken cases $N_{\mathit{broken}}(class_i)$, and the corresponding breakage rate $BR(class_i)$. If confidence intervals are ultimately computed, they should be attached to these rates rather than to isolated anecdotes. The key interpretive question is whether proof brittleness is concentrated in one pathological class or whether non-trivial breakage appears across multiple transformation families. Evidence for the latter would come from several classes exhibiting materially non-zero $BR(class_i)$ values and from those effects persisting across more than one seed-program stratum rather than being confined to tutorial-scale examples. When some classes remain low frequency in an early pilot or first benchmark release, the paper should preserve them in the manifest-level data but may merge or defer them in the main text tables if that yields a clearer presentation.

\subsection{Baseline Recovery Comparison}

We next evaluate baseline recovery on the broken subset induced by B1. In practice, the informative recovery quantities in early experiments are $RR(B2 \mid \mathit{broken})$ and, when enabled, $RR(B3 \mid \mathit{broken})$. B1 is still reported, but primarily as the mechanism that defines which transformed instances require recovery rather than as a genuine repair baseline. The central comparison is therefore not only which enabled recovery baseline attains the highest conditional recovery rate, but also how their recovered subsets overlap, which baselines are intentionally deferred and thus recorded as \texttt{not\_run}, and which benchmark items remain unrepaired after the enabled baseline set $\mathcal{B}_{\mathit{enabled}}$ has been applied. That remaining subset defines $Gap_{\mathit{residual}}(\mathcal{B}_{\mathit{enabled}})$. A strong version of the paper's empirical claim is that this residual gap remains substantial across multiple refactoring classes and seed strata even after nearby-version transfer and lightweight generation have been given a fair chance. If that pattern holds, then the benchmark shows not merely that breakage exists, but that simple recovery strategies leave meaningful nearby-version proof breakage unresolved.

\begin{table}[t]
\centering
\small
\caption{Planned breakage and recovery summary. The table is intended to show which refactoring classes are most brittle and how much of that breakage is removed by B1, B2, and B3 before the residual repair gap is measured.}
\label{tab:breakage-recovery}
\begin{tabular}{lcccccc}
\toprule
Class & Att. & Brk. & $BR(class_i)$ & $RR(B1)$ & $RR(B2)$ & $RR(B3)$ \\
\midrule
Rename & \tbdcell{$N^{\mathrm{rename}}_{\mathit{attempted}}$} & \tbdcell{$N^{\mathrm{rename}}_{\mathit{broken}}$} & \tbdcell{$BR(\mathrm{rename})$} & \tbdcell{$RR^{\mathrm{rename}}(B1)$} & \tbdcell{$RR^{\mathrm{rename}}(B2)$} & \tbdcell{$RR^{\mathrm{rename}}(B3)$} \\
Helper extraction & \tbdcell{$N^{\mathrm{extract}}_{\mathit{attempted}}$} & \tbdcell{$N^{\mathrm{extract}}_{\mathit{broken}}$} & \tbdcell{$BR(\mathrm{extract})$} & \tbdcell{$RR^{\mathrm{extract}}(B1)$} & \tbdcell{$RR^{\mathrm{extract}}(B2)$} & \tbdcell{$RR^{\mathrm{extract}}(B3)$} \\
Inlining & \tbdcell{$N^{\mathrm{inline}}_{\mathit{attempted}}$} & \tbdcell{$N^{\mathrm{inline}}_{\mathit{broken}}$} & \tbdcell{$BR(\mathrm{inline})$} & \tbdcell{$RR^{\mathrm{inline}}(B1)$} & \tbdcell{$RR^{\mathrm{inline}}(B2)$} & \tbdcell{$RR^{\mathrm{inline}}(B3)$} \\
Statement reordering & \tbdcell{$N^{\mathrm{reorder}}_{\mathit{attempted}}$} & \tbdcell{$N^{\mathrm{reorder}}_{\mathit{broken}}$} & \tbdcell{$BR(\mathrm{reorder})$} & \tbdcell{$RR^{\mathrm{reorder}}(B1)$} & \tbdcell{$RR^{\mathrm{reorder}}(B2)$} & \tbdcell{$RR^{\mathrm{reorder}}(B3)$} \\
Loop reshaping & \tbdcell{$N^{\mathrm{loop}}_{\mathit{attempted}}$} & \tbdcell{$N^{\mathrm{loop}}_{\mathit{broken}}$} & \tbdcell{$BR(\mathrm{loop})$} & \tbdcell{$RR^{\mathrm{loop}}(B1)$} & \tbdcell{$RR^{\mathrm{loop}}(B2)$} & \tbdcell{$RR^{\mathrm{loop}}(B3)$} \\
Contract factoring & \tbdcell{$N^{\mathrm{contract}}_{\mathit{attempted}}$} & \tbdcell{$N^{\mathrm{contract}}_{\mathit{broken}}$} & \tbdcell{$BR(\mathrm{contract})$} & \tbdcell{$RR^{\mathrm{contract}}(B1)$} & \tbdcell{$RR^{\mathrm{contract}}(B2)$} & \tbdcell{$RR^{\mathrm{contract}}(B3)$} \\
Ghost-code movement & \tbdcell{$N^{\mathrm{ghost}}_{\mathit{attempted}}$} & \tbdcell{$N^{\mathrm{ghost}}_{\mathit{broken}}$} & \tbdcell{$BR(\mathrm{ghost})$} & \tbdcell{$RR^{\mathrm{ghost}}(B1)$} & \tbdcell{$RR^{\mathrm{ghost}}(B2)$} & \tbdcell{$RR^{\mathrm{ghost}}(B3)$} \\
\midrule
Residual gap & \tbdcell{$N_{\mathit{broken}}$} & \tbdcell{$N_{\mathit{unrecovered}}$} & \tbdcell{$-$} & \multicolumn{3}{c}{\tbdcell{$Gap_{\mathit{residual}}$}} \\
\bottomrule
\end{tabular}
\end{table}


\subsection{Cost and Effort Proxies}

The recovery comparison should be read together with lightweight cost proxies. For each enabled baseline, we record verifier wall-clock time, and for B3 we additionally record retry count, prompt-token usage, completion-token usage, and total token-derived cost. These quantities are not direct measures of human effort, but they do help distinguish cheap robustness from expensive recovery. In particular, the paper should report whether improvements in $RR(B3 \mid \mathit{broken})$ come with large retry or token overhead relative to B1 and B2, and whether those costs vary meaningfully across refactoring classes or seed-program strata. If B3 is deferred in an early experiment tier, its result cells should appear explicitly as \texttt{not\_run} rather than being left blank; this keeps the distinction between absence of execution and failed recovery visible in both tables and prose.

\subsection{Failure-Mode Taxonomy and Residual Gap}

To move beyond aggregate success rates, we analyze the broken subset through the failure-mode labels introduced in Section~3. Let $L_{\mathit{assert}}$, $L_{\mathit{inv}}$, $L_{\mathit{lemma}}$, $L_{\mathit{trigger}}$, $L_{\mathit{ghost}}$, and $L_{\mathit{tool}}$ denote the counts associated with the six label categories. These correspond respectively to broken assertion flow, weakened loop invariants, brittle lemma or helper-application structure, trigger or quantifier sensitivity, ghost-code or proof-hint displacement, and unexplained or tool-sensitive failure. This view helps distinguish which kinds of nearby-version proof breakage are partially addressed by surface transfer and which persist into $Gap_{\mathit{residual}}(\mathcal{B}_{\mathit{enabled}})$ even after the enabled recovery baselines have been applied. In particular, a residual concentration in lemma-structure, invariant, or trigger-sensitive cases would motivate stronger symbolic repair in later work more convincingly than a residual dominated by simple renaming failures. If some labels appear only once or twice in an early experiment tier, they should remain visible in the raw results but may be summarized more coarsely in the main narrative to avoid overweighting pilot-scale noise.

\begin{table}[t]
\centering
\caption{Planned failure-mode distribution. The table is intended to show which proof-relevant failure categories dominate the broken subset and which categories continue to populate the residual repair gap after baseline recovery.}
\label{tab:failure-modes}
\begin{tabular}{lrr}
\toprule
Failure-mode label & Broken subset & Residual gap subset \\
\midrule
Broken assertion flow & \tbdcell{$L_{\mathit{assert}}$} & \tbdcell{$L^{\mathrm{res}}_{\mathit{assert}}$} \\
Weakened loop invariants & \tbdcell{$L_{\mathit{inv}}$} & \tbdcell{$L^{\mathrm{res}}_{\mathit{inv}}$} \\
Brittle lemma / helper structure & \tbdcell{$L_{\mathit{lemma}}$} & \tbdcell{$L^{\mathrm{res}}_{\mathit{lemma}}$} \\
Trigger / quantifier sensitivity & \tbdcell{$L_{\mathit{trigger}}$} & \tbdcell{$L^{\mathrm{res}}_{\mathit{trigger}}$} \\
Ghost-code / proof-hint displacement & \tbdcell{$L_{\mathit{ghost}}$} & \tbdcell{$L^{\mathrm{res}}_{\mathit{ghost}}$} \\
Unexplained / tool-sensitive failure & \tbdcell{$L_{\mathit{tool}}$} & \tbdcell{$L^{\mathrm{res}}_{\mathit{tool}}$} \\
\bottomrule
\end{tabular}
\end{table}


At the same time, unrepaired cases should not automatically be interpreted as semantically deep failures. Some may still reflect solver brittleness, hidden context sensitivity, or benchmark-local effects that fall outside the expressive reach of B1--B3 without requiring fundamentally new semantics-aware reasoning. The taxonomy is therefore used here as an analysis device for interpreting the residual repair gap, not as a definitive theory of proof failure.

\subsection{Takeaways}

Taken together, the results should make one claim legible even before exact numbers are inserted: Dataset~A exposes nearby-version proof breakage as a measurable maintenance problem, that problem appears across multiple semantics-preserving refactoring classes rather than in a single corner case, and the enabled recovery baselines still leave a non-trivial residual repair gap. In other words, the benchmark is designed to show not just that verified code can break under evolution, but that nearby-version proof maintenance remains a structured empirical problem after the easiest fixes have already been tried. That is the empirical substrate this paper aims to establish so that later repair work can be evaluated against it rather than inferred from isolated examples.
